\documentclass[11pt]{article}

\usepackage[margin=25mm]{geometry}

\usepackage{amsmath,amssymb,fontspec,unicode-math}

\IfFontExistsTF{DejaVu Serif}{\setmainfont{DejaVu Serif}}{\setmainfont{Latin Modern Roman}}

\IfFontExistsTF{Latin Modern Math}{\setmathfont{Latin Modern Math}}{\setmathfont{DejaVu Math TeX Gyre}}

\usepackage{xurl,hyperref}

\setlength{\parindent}{0pt}\setlength{\parskip}{6pt}

\begin{document}

\title{TP-MDDN: Task-Preferenced Multi-Demand-Driven Navigation with Autonomous Decision-Making}

\author{ELVIA preliminary AI analysis}\date{21 October 2026}\maketitle

Preliminary assessment: Minor.\par

Independent scientific review remains pending. Author adoption has not been requested.

Under the natural finite-navigable-grid normalization, zero obstacle preference conflates a safe minimum with threshold exclusion. Constant distance ranges also leave the stated formulas undefined. The impact on actual navigation remains conditional on implementation conventions.

\section*{Review scope}

Claim-focused reading of the original paper and relevant author-released source, with direct mathematical and static-source checks. No author program, historical verifier, model training or published benchmark was executed.

\section*{Problem 1: A normalized minimum does not encode the raw obstacle threshold}

Source: Section 3.1.3, Calculation of Affordance Map and Eq.(3), PDF p.7; Section 4.1 threshold, p.7

Assessment: supported; impact: minor

\subsection*{Source claim (paraphrase)}

For threshold-admissible navigable points the obstacle preference is min–max normalized, but Eq.(3) uses its zero value as the final exclusion mask. Target-distance normalization also divides by its range without a stated tie rule.

\subsection*{Analysis}

Let the already processed obstacle cloud be \(O=\{(0,0)\}\), target \(T=(1,0)\), candidates \(n_1=(1/2,0)\), \(n_2=(3/4,0)\), and threshold \(\tau_o=1/4\). Both are safe under the stated raw-distance threshold. Their obstacle distances are \((1/2,3/4)\), hence normalized preferences \((0,1)\); their inverse target preferences are \((0,1)\). Eq.(3) returns \((0,1)\), whereas a raw-distance mask followed by the stated clipping returns \((0.1,1)\). Every safe minimum is similarly excluded when \(d_{\max}>d_{\min}\ge\tau_o\).

Valid geometry also gives undefined ranges. Candidates \((\pm1/2,0)\) around the same obstacle have constant safe obstacle distance and yield \(0/0\). Alternatively, the original candidates with target \((5/8,0)\) have equal target distances \(1/8\), so semantic normalization yields \(0/0\).

\subsection*{Finding and scope}

The literal formulas need a normalization and tie convention. If excluding safe minima is intentional, it is an additional preference rule; normalization over a larger set containing obstacles could remove the first example. No route or success-rate change is established.

\subsection*{Proposed correction}

Separate the mask \(s(n)=\mathbf1\{d_o(n)\ge\tau_o\}\) from preference. Let \(q\) be inverse target normalization with an explicit common value when distances tie, and set \(\widehat a_f=s\,\operatorname{clip}(q,0.1,1)\). The positive clipped factor makes zero equivalent to raw threshold exclusion. This restores only that property, not collision-free planning.

A proposed correction is described; author adoption is not established.

\section*{Limits}

Preliminary AI assessment; independent human scientific adjudication remains pending.\par

Analytic examples and release discrepancies do not remeasure the paper’s reported experiments.\par

The actual navigation implementation and any unpublished normalization or constant-range guards were not inspected.\par

\section*{Source identity}

arXiv:2511.17225v1 (current author version confirmed from official metadata)

\url{https://arxiv.org/pdf/2511.17225v1}

Reviewed PDF SHA-256: \nolinkurl{70df9aee25cd8defbc291bc1ff4ee1e327c356f2caeb00a877c2553a3745fe45}

The SHA-256 identifies the exact retained PDF inspected in this rapid source/argument check. Original supplied report bytes are retained locally. No unambiguous original PDF retrieval timestamp was recovered; file modification time was not substituted. Current official arXiv metadata checked at 2026-10-21T05:36:46.042465+00:00 identifies 2511.17225v1; implicated passages were checked in the current identified material.

\end{document}
